\documentclass[10pt]{scrreprt}
\usepackage[a4paper, top=30mm, left=25mm, right=25mm, bottom=30mm]{geometry}
\usepackage[utf8]{inputenc}

\usepackage[Bjornstrup]{fncychap}
\usepackage{pdfpages}
\usepackage{ngerman}
\usepackage{graphicx}
\usepackage{epstopdf}
\usepackage{etoolbox}
\usepackage{wrapfig}
\usepackage{enumitem}
\usepackage{tabu}
\usepackage{amsmath}
\usepackage{url}
\usepackage{numprint}
\usepackage{longtable}
\usepackage{caption}
 \usepackage{rotating}
\usepackage{array}
\usepackage{listings}
\usepackage{amssymb}
\usepackage{multirow}
\usepackage{textcomp}
\usepackage{calc}
\usepackage{ulem}
\usepackage[table]{xcolor}
%\usepackage{hyperref}
%\hypersetup{%
%  colorlinks=false,%
%  linkcolor=blue,%
%  citecolor=blue,%
%  unicode%
%}
\usepackage{mathptmx}
\usepackage{courier}
\usepackage{amssymb}

\renewcommand{\em}{\sffamily\bfseries}
\renewcommand{\bf}{\sffamily\bfseries}
\newcommand{\hyperlink}[2]{#1}

\makeatletter
\patchcmd{\@makechapterhead}{\vspace*{50\p@}}{\vspace*{-20\p@}}{}{}
\patchcmd{\@makeschapterhead}{\vspace*{50\p@}}{\vspace*{7\p@}}{}{}
\patchcmd{\DOTIS}{\vskip 40\p@}{\vskip -12\p@} 
\makeatother
  
\captionsetup[figure]{labelfont={sf,bf},textfont={sf}}
\deffootnotemark{[\thefootnotemark]}
\deffootnote{1.5em}{1em}{[\thefootnotemark] }
\setlength{\parindent}{0pt}
\renewcommand{\labelitemi}{ \raisebox{0.3ex}{\small$\blacktriangleright$} }
\renewcommand{\labelitemii}{ \raisebox{0.3ex}{\small$\triangleright$} }
\lstset{language=Java}

\newcommand{\sfbf}[1]{\textbf{\sffamily #1}}
\newcommand{\sfit}[1]{\textit{\sffamily #1}}
\newcommand{\W}{\sfbf{W}}
\newcommand{\ziel}[1]{{\fontsize{9.5}{11}\textsf{/#1/}}}
\newcommand{\ziellabel}{Z}
\newcommand{\muss}{\renewcommand{\labelenumi}{\textbf{\ziel{\ziellabel\numprint{\theenumi}0}}}}
\newcommand{\wunsch}{\renewcommand{\labelenumi}{\textbf{\ziel{\ziellabel\numprint{\theenumi}0W}}}}
\newcommand{\JoglEarth}{\raisebox{-1.2mm}{\includegraphics[scale=0.33]{Logo-Text.eps}} }
\newcommand{\textref}[1]{\mbox{\raisebox{0.1ex}{\small$\rightarrow$ }\textit{#1}}}

\newenvironment{details}[1][6pt]{%
  \parskip#1 \parindent6mm \raggedright%
  \def\item{\par\ignorespaces\hangindent=5mm \hangafter1}}{%
  \par\ignorespaces} 
  
 
\begin{document}

\thispagestyle{empty}
\sffamily
 
\title{Validierung}

\begin{figure}
\begin{flushright}
	\includegraphics[scale=0.4]{uniLogo.eps}
\vspace{2.0 cm}
\end{flushright}
\end{figure}

\begin{center}
\vspace{2.0 cm}
{\LARGE SEP – Wintersemester 2013/14}

\vspace{1.0 cm}
\textbf{{\Huge Validierungsbericht}}

\vspace{0.4 cm}
\begin{figure}[!htb]
\begin{center}
	%\includegraphics[scale=1.0]{projektLogo.eps}
	\includegraphics[scale=1.5]{Logo-Print.eps}
\end{center}
\end{figure}

\vspace{0.2 cm}
\textbf{{\huge OpenStreetMap: Die Welt in 3D}}

\vspace{1.5 cm}
17.1.2014

\vspace{0.5 cm}
Version: 1.0

\vspace{1.5 cm}
{\Large Projektbetreuer: Peter Barth}

\vspace{1.5 cm}
\begin{tabular}{|c|c|c|}
\hline 
\rule[-1ex]{0pt}{4ex} \textbf{Phase} & \textbf{Verantwortlicher} & \textbf{E-Mail Adresse} \\ 
\hline  \hline
\rule[-1ex]{0pt}{4ex} Pflichtenheft & Gabriele Haas & haasgab@fim.uni-passau.de \\ 
\hline  \hline
\rule[-1ex]{0pt}{4ex} Entwurf & Thomas Eder & ederthom@fim.uni-passau.de \\ 
\hline  \hline
\rule[-1ex]{0pt}{4ex} Spezifikation & Christof Blauberger & blauberg@fim.uni-passau.de \\ 
\hline  \hline
\rule[-1ex]{0pt}{4ex} Implementierung & Fabian Knorr & knorrfab@fim.uni-passau.de \\ 
\hline \hline 
\rule[-1ex]{0pt}{4ex} Testing & Constantin Wenger & wengerco@fim.uni-passau.de \\ 
\hline  \hline
\rule[-1ex]{0pt}{4ex} Präsentation & Sebastian Reichl & reichlse@fim.uni-passau.de \\ 
\hline 
\end{tabular}

\end{center}


\pagebreak
\rmfamily
\tableofcontents

\chapter{Einleitung}

\chapter{Testfälle im Pflichtenheft}
\section{Vor der Validierung}
\subsection{Anwendungsstart}
/T050/: Seitenleiste ist zu klein nach dem wieder einblenden.
\subsection{Globusansicht}
/T150/: Kacheln dauern sehr lange zum laden, da welche davon 404 zurückhaben und für Minuten versucht wurde diese zu laden.
\subsection{Navigation auf dem Globus}
/T230W/: Höhenprofil scheint nicht aktiviert zu werden
\subsection{Overlays}
Keine Overlays tauchen auf
\subsection{Locale Suche}
/T570/ Ändern des Zoomlevels resultierte in OpenGL out of Memory Fehler
\subsection{Globale Suche}
/T650W/ Suchergebnisse tauchen nie auf
\subsection{Punkte Markieren}
/T660W/ GL Out of Memory Fehler
\subsection{Caching, Diverses}
/T800/ erstmal keine Textur
/T820/ Ordner Größe verringert sich nicht
/T850/ Passiert nicht

\section{Nach der Validierung}

\chapter{Bugs}
\section{Gefundene Bugs}
Die Nummerierung wurde aus dem Bugtracker von GitHub entnommen.
Nicht in der Liste verzeichnet sind Verbesserungen und Duplikate.
\begin{longtable}{|c|p{5.2cm}|p{8.2cm}|}
\hline
Nummer & Titel & Beschreibung \\
\hline
\hline
\#1 & Sternenhimmel funktioniert unter Windows 8.1 nicht & \\
\hline
\#2 & Zu viele Nominatim-Anfragen & Es wird für jeden Pixel eine Anfrage gestellt \\
\hline
\#3 & Settings werden nicht korrekt in das UI geladen & \\
\hline
\#4 & Detaillevel-Einstellung stimmt nicht mit dem UI überein & Es wird 'Niedrig' angezeigt
 bwohl 'Mittel' eingestellt ist \\
\hline
\#5 & Letzter Frame wird nicht angestoßen & Bei der Umstellung zu Globus oder Einblendung von Overlays wird der Rendervorgang nicht angestoßen\\
\hline
\#6 & Kugel wird bei kleinem Zoomlevel zu niedrig aufgelöst & \\
\hline
\#7 & Sichtbarkeitsberechnung funktioniert für die Pole nicht & Pole sind praktisch immer unsichtbar; Zentrieren einer Polkachel lässt Ebene/Kugel verschwinden \\
\hline
\#8 & Kugel zeigt bei bestimmtem Zoomlevel Artefakte & Zoomlevel 3 oder 4 führt dazu, dass am 60. Breitengrad Artefakte durch die Längengradreduktion des 'SphereTessellators' auftreten \\
\hline
\#9,\#10 & Layout der Oberfläche wird nicht korrekt angezeigt & Bei gewissen Fenstermanagern wird die Oberfläche nicht korrekt angezeigt; es wird ein grau/blauer Balken angezeigt der stört. Deshalb wird ein Stück der Oberfläche abgeschnitten \\
\hline
\#13 & Implementierung der Maßstabsanzeige zeigt Schwächen & \\
\hline
\#16 & Bestimmte PDF-Reader öffnen Handbuch inkorrekt & Es ist z.B. der Adobe Reader 9 unter Linux betroffen \\
\hline
\#17 & Sprachänderung inkonsistent & Bei der Änderung der Sprache von Deutsch nach Englisch wurden z.B. 'Längengrad' und 'Breitengrad' nicht übersetzt\\
\hline
\#18 & Zoombalken stimmt noch nicht mit Zoom bei OSM überein & Zoombalken 100\% sollte 'OSM' Zoomstufe 18 sein. Momentan Zoomstufe 14 maximal möglich\\
\hline
\#19 & Höhenprofil funktioniert noch nicht & Kacheln werden manchmal geladen, aber zu viele und sie werden nicht verarbeitet \\
\hline
\#20 & Kippen der Kamera mit Tastatureingaben zum Teil seltsam & Pos1, Ende, Bild Aus, Bild Ab macht seltsame Bewegungen (Zu ruckartig) \\
\hline
\#21 & Noch nicht implementierte Tastatureingaben & Numpad 0 zum Zurücksetzen des Kippens und F1 zur Anzeige des Handbuchs noch nicht implementiert \\
\hline
\#22 & Details zu Punkt noch nicht angezeigt & Detailfeld bleibt leer \\
\hline
\#23 & Knopf zum Löschen der Benutzermarkierungen fehlt & Löschen der Benutzermarkierungen noch nicht implementiert \\
\hline
\#24 & Markierung für Bildmittelpunkt nicht am korrekten Platz angezeigt & Overlays wurden an falsche Stellen gezeichnet \\
\hline
\#26 & Nach Auswahl des Kartentyps bewirkt Pfeil auf und Pfeil ab eine erneute Änderung des Typs & Statt die Karte zu bewegen wird durch die Kartentypen navigiert \\
\hline
\#27 & Exception im MainWindow bei initialem Programmstart & Ausgelöst durch fehlende Collection 'userLocations' \\
\hline
\#28 & Exception bei Nominatimanfragen & Nach Anwahl eines Ergebnisses entsteht eine Exception bei der zweiten Suche \\
\hline
\#30 & Aus- und Einblenden der Seitenleiste zerstört das Layout der Oberfläche & Nach Aus- und wieder Einblenden der Seitenleiste wird die Seitenleiste zu schmal angezeigt \\
\hline
\#31 & Kippen funktioniert nicht konsistent mit der Drehung & \\
\hline
\#32 & Verschieben ist nicht konsistent mit der Kartenoberfläche & Punkt unter der Maus sollte beim Verschieben unter selbiger bleiben \\
\hline
\#33 & Drehen über 180° E funktioniert mit Tastatur nicht & Navigation über Karte nicht mehr möglich, nachdem über 180° E geschoben wurde \\
\hline
\#34 & Liste der Auswahlmöglichkeiten verschiebt Labels bei langen Suchergebnissen & Die Labels 'In der Nähe', 'Weltweit' und der Suchknopf wurden verschoben \\
\hline
\#43 & GL Out of Memory Error tritt selten und zufällig auf & \\
\hline
\#44 & NumberFormatException bei 'GeoCoordinates' wird nicht gefangen & Beim Parsen der Koordinaten wird nicht auf Gültigkeit überprüft \\
\hline
\#45 & Position des Mittelpuktes bei Kartenansicht im Zoomlevel 0/1/2/3 stimmt nicht & \\
\hline
\#47 & 'ProgressManager' wird nach vollständigem Laden nicht zurückgesetzt & Ist 100\% einmal erreicht, sollte eine neue Anfrage den Fortschrittsbalken auf 0\% setzen \\
\hline
\#50 & Overpass-Anfragen dauern zu lange & Bei Einschalten der Overlays dauern die Anfragen zu Overpass ungewöhnlich lange \\
\hline
\#56 & Bildschirmecken werden bei der Sichtbarkeitsberechnung falsch behandelt & Es entstehen schwarze Löcher, obwohl die Kacheln sichtbar sein sollten \\
\hline
\#57 & Sichtbarkeitsberechnung der Ebene fehlerhaft. Es werden generell alle Kacheln geladen & \\
\hline
\#58 & Lokale Suche funktioniert nicht & Ergebnisse erscheinen kurz im Auswahlfenster, verschwinden dann aber wieder \\
\hline
\end{longtable}
\section{Verbesserungsvorschläge}
\begin{longtable}{|c|p{5.2cm}|p{8.2cm}|}
\hline
Nummer & Titel & Beschreibung \\
\hline
\hline
\#11 & Mausrad-Zoomen sollte auf Mausposition statt Mittelpunkt arbeiten & Zur Erhöhung der Benutzerfreundlichkeit sollte auf die Mausposition gezoomt werden \\
\hline
\#12 & In der Kugelansicht sollte das Zoomlevel anders berechnet werden & Derzeitige Berechnung führt zu Unregelmäßigkeiten in Grenzregionen \\
\hline
\#13 & Ausgrauen von Längen- und Breitengradeingabe im Sonnensystemmodus & Eingabe des Benutzers in dieser Ansicht eventuell nicht gewünscht \\
\hline
\#15 & Mond im Sonnensystemmodus entzerren & Er scheint etwas in der Form verunstaltet zu sein \\
\hline
\#25 & Die Einheit der Cachegrößen sollte angezeigt werden & Momentan sind keine Einheiten angegeben \\  
\hline
\#26 & Hineinzoomen beim Sonnensystemmodus sollte in die Globusansicht wechseln & Durch die niedrige Auflösung der mitgelieferten Textur sollte bei hohem Zoomlevel in die Globusansicht gewechselt werden \\
\hline
\#36 & Fehlende OpenGL-Unterstützung sollte eine Meldung geben & Falls ein System OpenGl 2.0 nicht unterstützt, sollte eine Meldung erscheinen \\
\hline
\#40 & Anzeige-Tab sollte eventuell Ansichts-Tab heißen & \\
\hline
\#41 & \JoglEarth bleibt hängen, sollte es eine Kachel nicht geben & \\
\hline
\#48 & Erdachse einzeichnen & Man könnte eine Erdachse einzeichnen, die ähnlich einem Kompass eingefärbt ist \\
\hline
\#49 & Schrittgröße im Spinner für Cachegröße &
Sollte in Zehnerschritten erfolgen \\
\hline
\#51 & Laden der Städtenamen bei niedrigen Zoomleveln & Anfragen bei Overpass dauern für große Bereiche sehr lange \\
\hline
\#52 & Reaktion auf HTTP-Fehler bei Suchanfrage & Bei Unerreichbarkeit von Servern oder Timeouts sollte noch eine Anfrage gesendet werden und dann ein leeres Suchergebnis geliefert werden, sollte das Problem weiterhin bestehen \\
\hline
\#53 & Globus wird niedriger Detailstufe zu einer Ebene & Es wird versucht aus 4 Punkten eine Kugel zu zeichnen \\
\hline
\#55 & Ausgewählte Overlays könnten gespeichert werden & Die Anwahl von Overlays könnten eventuell über das Beenden des Programms hinaus gespeichert werden um dann beim Start geladen zu werden \\
\hline
\end{longtable}

\chapter{Code Coverage}
\section{Pflichtenheft-Tests}
\section{JUnit-Test}
\section{Test per Hand}
\chapter{Skalierungs- und Belastungstest}
\chapter{Profiling}
\end{document}